\chapter{Conclusion and Future Work}
\label{Chap6}

\section{Conclusion}

This dissertation began with a simple expectation: if several
vision--language models received the same image, moral taxonomy, and
instructions, their agreement---especially confident agreement---might indicate
a shared moral reading. The findings complicate that expectation. Rather than
asking which model produced the correct label, the study examined how moral
interpretations differed across three model configurations, participant
responses, repeated runs, and selected changes to visual context.

RQ1 and RQ2 distinguish two aspects of visual moral attribution: whether an
image was classified as morally relevant and, if so, how its moral meaning was
framed. The configurations crossed the first boundary at
different rates on the fixed corpus, but the uneven stimulus set means that
these output margins cannot establish general model preferences. Agreement
also increased when the ten moral labels were collapsed into moral content
versus \textit{None}. That increase did not resolve the original
interpretations; it removed distinctions among them. Some disagreements
contrasted vulnerability with suffering within Care/Harm, whereas others
placed the same scene in different foundations. \textbf{The models agreed more
readily that a scene had moral significance than on what that significance
meant.}

Repeated runs revealed a second distinction. Within the later repeat block,
each configuration usually reproduced its own labels more often than it matched
another model. Yet this short-run consistency did not always reproduce the
earlier primary output. The contrast was clearest for GPT-5.6, although the two
blocks differed in date and settings and cannot support a general temporal
ranking. \textbf{A model can therefore be repeatable within one evaluation
window without preserving the same interpretation across evaluation windows.}

RQ3 supplied the most important human comparison. Participants did not form a
single answer key: on eight of the 16 survey images, even the most-endorsed
category was selected by fewer than half of those who saw the image. Within
this purposively selected subset, GPT-5.6 nevertheless showed the highest
descriptive correspondence with participant selections on all four summary
measures. This was a limited advantage. All three models returned the same
label on ten images, so the ordering was produced entirely by the remaining six
model-disagreement cases. Model consensus and reported confidence also failed
to guarantee strong participant support. Thus, GPT-5.6 corresponded most
closely with these participants on these images; the result does not establish
that it is generally closest to human moral judgement.

RQ4 then showed how a moral reading could change. In the selected cases,
obscuring a potentially relevant cue often preserved the original label,
although confidence sometimes fell. By contrast, replacing distressed expressions with positive ones or changing
the background context could reorganise the classification and expose
different alternative readings across models. This
pattern points to relationships among visible cues rather than a single
cue-to-label trigger, but the small, imperfectly isolated edits cannot identify
a causal visual mechanism. The supplementary prompt check produced an equally
narrow result: the structured and cue-guided prompts caused only small, mixed
changes. GPT-5.6 retained the highest descriptive correspondence under all
three conditions, but cue guidance did not consistently improve participant
correspondence.

Taken together, the results support one central conclusion: \textbf{the models
did not simply retrieve a moral category already contained in an image; they
constructed a moral reading by selecting which visible evidence and contextual
inferences mattered.} Agreement about moral presence could conceal disagreement
about meaning, local repeatability could conceal weaker continuity, and model
consensus could conceal a different distribution of participant
interpretations. Agreement, participant correspondence, repetition, and
confidence are therefore distinct signals, not one measure.

\section{Scope and Limitations}

These conclusions apply to the recorded configurations and study materials.
The 127-image corpus was constructed rather than sampled from a defined image
population, and its researcher-assigned indexing categories were highly uneven
and not independently validated. Of the 127 corpus images, 76 were generated
using Google Gemini, so generator-specific stylistic regularities or visual
artefacts may have influenced the model outputs; the study did not separately
estimate any such influence. The primary analysis retained one output for
each model--image pair. The repeat subset was selected from initial agreement
patterns, and the primary and repeat blocks used non-identical settings.
Consequently, neither component estimates general model preferences or
population-level stability.

The participant comparison used 16 purposively selected images, 56
self-selected respondents, and no demographic variables. Participants could
select several categories under shorter definitions, whereas each model had to
return one label under more detailed instructions. Participant endorsement was
therefore a measure of correspondence under this survey design, not accuracy
under matched conditions. The image edits were few and did not hold every
non-target feature constant, while the prompt check reused the same small
survey subset. These limitations do not change the observed descriptive
patterns, but they restrict how far those patterns can be generalised. The
absence of a clear prompt improvement also does not show that a better prompt
could not be developed.


\section{Future Work}
\label{sec:future_work}

The first priority is a larger and more balanced image corpus. A planned
diagnostic set should provide adequate coverage of moral and non-moral scenes,
all five foundations, both virtue and vice poles, and deliberately ambiguous
cases. A separate naturalistic sample should then test whether findings from
the balanced set remain visible when category frequencies are not controlled.
Keeping development and evaluation images separate would make it easier to
distinguish a recurring model pattern from a feature of the particular images
used to discover it.

More human annotation is equally important. Future studies should obtain more
independent judgements for every image and collect short reasons or ranked
interpretations alongside category selections. Relevant demographic variables
would allow planned analyses of cultural variation
\cite{atari2023weird,yadav2025beyond}. Models and participants should also
receive more closely matched definitions and response options, including a
clear distinction between no moral content and insufficient visible evidence.
The aim would not be to manufacture one universal human label, but to estimate
which readings are widely shared, contested, or associated with particular
groups and interpretations.

The image-edit findings provide concrete hypotheses for both controlled visual
tests and prompt engineering. Future counterfactual pairs should be created
before inference, change one intended cue at a time, and include matched control
edits. Repeated evaluation of both images would help separate sensitivity to
the intended cue from ordinary between-call variation
\cite{liu2026moralbackbone}. Prompt rules could then be developed on a larger,
balanced development set by examining recurring differences between visible
evidence and inferred background stories. The prompt should be frozen before
testing on separate held-out images. The present P1 and P2 conditions did not
produce a consistent improvement; a larger dataset would not guarantee a
better prompt, but it would provide a stronger basis for testing whether the
cue-related hypotheses generalise. Prompt engineering should therefore remain
a robustness check rather than a post-hoc search for preferred answers.

Finally, future evaluations should continue to separate short-run
repeatability from longer-term continuity. More repetitions, broader image
coverage, and date-stamped reruns of the same frozen requests would help locate
variation associated with the image, prompt, model configuration, or
evaluation window. Retaining provider settings, returned model identifiers,
failures, and raw outputs would make later changes in hosted systems easier to
detect \cite{chen2024changing}.

Together, these extensions would move the evaluation beyond asking whether
models return the same label. The more informative question is which evidence
makes a visual scene morally meaningful to different models and different
people, and under what conditions that interpretation remains stable.

\section{Implications for the Sustainable Development Goals}
\label{sec:sdg_implications}

The findings have their clearest methodological relevance to SDG~16,
particularly Targets~16.6 and~16.7. Where LVLMs might inform institutional
judgements, evaluations intended to support accountable, transparent, and
inclusive use should keep interpretive disagreement visible.
\textbf{Model consensus or model-reported confidence should therefore not be
treated as evidence of a uniquely correct moral interpretation.}

The connection to SDG~10 is secondary and indirect. Variation among models and
participants raises the question of whose interpretations are represented when
a system assigns a single label, but the study does not demonstrate unequal
outcomes or effects on any group. These implications concern evaluation
practice; the dissertation does not measure progress towards either SDG.
